% ============================================================
%  ch3_II  第二章 证明的大线条（原书 pp. 12--19，即 PDF 第 13--20 页）
% ============================================================

\section{证明的大线条}
\label{sec:II}
\origpage{12}

\par\medskip
\noindent\textbf{边界为 $r$-可忽略的紧化}

\par\medskip
\noindent\textbf{$\bullet$ 无水平的情形：$N = \emptyset$}

对每个截断多边形 $p : [0, r] \to \RR_+$（足够凸），存在 $\Cht^{r,\ol{p}\le p}/a^{\ZZ}$ 的一个紧化

\[ \ol{\Cht}^{r,\ol{p}\le p}/a^{\ZZ} \cong \coprod_{1 \le d \le r|\deg(a)|} \ol{\Cht}^{r,d,\ol{p}\le p} \]

使得：

\begin{itemize}
\item $\ol{\Cht}^{r,d,\ol{p}\le p}$ 在 $X \times X$ 上紧合且光滑；
\item 它在
\[ X \times X \times \Aa^{r-1}/\mathbb{G}_m^{r-1} \]
 上也光滑，其中边界 = 正规交叉除子（记作 $\mathrm{bord}$）；
\item 它的边界分层由划分
\[ \ul{r} = (r = r_1 + \cdots + r_k) \rightsquigarrow \text{记为 } \Cht^{r,d,\ol{p}\le p}_{\ul{r}} \]
 编号，并且本质上形如
\[ \Cht^{r_1,d_1,\ol{p}\le p_1} \times_X \cdots \times_X \Cht^{r_k,d_k,\ol{p}\le p_k}; \]
\origpage{13}
 它们带有到下面对象的态射：
\[
\begin{array}{ccc}
X \times & X^{k-1} \times & X\\
\uparrow & \uparrow & \uparrow\\
\text{极点} & \text{“退化子”} & \text{零点}
\end{array}
\]
\end{itemize}

\begin{proposition*}
在归纳假设下，边界分层
\[ \Cht^{r,d,\ol{p}\le p}_{\ul{r}} \]
的上同调（$\ell$-进、紧支集）是 $r$-可忽略的。
\end{proposition*}

\par\medskip
\noindent\textbf{证明.}\ 分两步计算上同调：
\[
\begin{array}{c}
\Cht^{r_1,d_1,\ol{p}\le p_1} \times_X \cdots \times_X \Cht^{r_k,d_k,\ol{p}\le p_k}\\[2pt]
\Big\downarrow\ \text{在 } X \times X^{k-1} \times X \text{ 上的上同调}\\[2pt]
(q_0^*\sigma_1 \otimes q_1^*\sigma'_1) \otimes \cdots \otimes (q_{k-1}^*\sigma_k \otimes q_k^*\sigma'_k) \quad (\dim \sigma_i, \sigma'_i < r,\ \forall i)\\[2pt]
\Big\downarrow\ \text{在 } X \times X \text{ 上的上同调}\\[2pt]
(\sigma_1 \otimes \sigma'_k) \otimes \lambda \quad (\lambda \text{ 为 } \widehat{\ZZ} = \Gal(\overline{\mathbb{F}}_q) \text{ 的特征})
\end{array}
\]

\par\medskip
\noindent\textbf{$\bullet$ 带水平 $N \neq 0$ 的情形：}

$\Cht^{r,\ol{p}\le p}_N/a^{\ZZ}$ 是 $\Cht^{r,\ol{p}\le p}/a^{\ZZ} \times_{X\times X} (X - N) \times (X - N)$ 的有限 étale Galois 覆盖（群为 $\GL_r(\cO_N)$）。设

\[ \ol{\Cht}^{r,\ol{p}\le p}_N/a^{\ZZ} = \coprod_{1 \le d \le r|\deg(a)|} \ol{\Cht}^{r,d,\ol{p}\le p}_N/a^{\ZZ} \]

为 $\Cht^{r,\ol{p}\le p}/a^{\ZZ}\times_{X\times X}(X - N)\times(X - N)$ 在 $\Cht^{r,\ol{p}\le p}_N/a^{\ZZ}$ 中的归一化。它在 $(X - N)\times(X - N)$ 上紧合，但\emph{并不光滑}。设

\[ \ol{\Cht}^{r,\ol{p}\le p'}_N/a^{\ZZ} = \coprod_{1 \le d \le r|\deg(a)|} \ol{\Cht}^{r,d,\ol{p}\le p'}_N \]

为由如下条件定义的开集：不仅极点 $\infty$ 与零点 $0$，而且\emph{退化子}都避开 $N$。

\begin{theorem*}[稳定性定理]
\emph{Hecke 对应 $f \in \cH^r_N$ 稳定开集 $\ol{\Cht}^{r,\ol{p}\le p}_N/a^{\ZZ}$。}
\end{theorem*}

\begin{proposition*}
\upshape (i) \emph{$\ol{\Cht}^{r,\ol{p}\le p}_N/a^{\ZZ}$ 在 $X \times X$ 上光滑，甚至在}
\[ X \times X \times \Aa^{r-1}/\mathbb{G}_m^{r-1} \]
\emph{上光滑；其边界为正规交叉除子（DCN），分层由划分 $\ul{r} = (r = r_1 + \cdots + r_k)$ 编号。}
\upshape (ii) \emph{若在归纳假设下，其边界分层的上同调是 $r$-可忽略的。}
\end{proposition*}

\origpage{14}

\begin{corollary*}
每个 $f \in \cH^r_N$ 都作用在 $\ol{\Cht}^{r,\ol{p}\le p}_N/a^{\ZZ}$ 的上同调上。
\end{corollary*}

记：

$\pi_1^{(X-N)\times(X-N)}$ 为 $(X - N) \times (X - N)$ 的基本群，

$W_{(X-N)\times(X-N)} = \pi_1^{(X-N)\times(X-N)} \times_{\widehat{\ZZ}} \ZZ$ 为其 Weil 群。

\begin{lemma*}
存在 $W_{(X-N)\times(X-N)}$ 的不可约表示的、系数在 $\QQ$ 中的一个形式线性组合 $H^*_{N,\mathrm{cusp}}$，使得：
\upshape (i) \emph{差}
\[ H^*_{N,\mathrm{cusp}} - \frac{1}{r!}\sum_{k=1}^{r!} (\Frob_X^k \times \mathrm{Id}_X)^*\, H^*_c(\ol{\Cht}^{r,\ol{p}\le p}_N/a^{\ZZ}) \]
\emph{的所有分量都是 $r$-可忽略的（其中 $H^*_c(\cdot) = \sum_\nu (-1)^\nu H^\nu_c(\cdot)^{ss}$ 是 $X \times X$ 泛点之上的“交错”上同调）。}
\upshape (ii) \emph{对每个 $x \in |X \times X|$，在其上 $\infty \neq 0 \in |X - N|$，并对 $\deg(x)$ 的每个倍数 $s = \deg(\infty)s' = \deg(0)u'$，}
\[ \Tr_{H^*_{N,\mathrm{cusp}}}\big(\Frob_x^{-\frac{s}{\deg(x)}}\big) = q^{(r-1)s} \sum_{\substack{\pi \in \{\pi\}^N_r\\ \chi_\pi(a) = 1}} \big(z_1(\pi_\infty)^{-s'} + \cdots + z_r(\pi_\infty)^{-s'}\big)\big(z_1(\pi_0)^{u'} + \cdots + z_r(\pi_0)^{u'}\big). \]
\end{lemma*}

\par\medskip
\noindent\textbf{所需的材料：}

\begin{itemize}
\item 不动点计数公式；
\item 对应 $\pi' \mapsto \sigma_{\pi'}$ 在秩 $< r$ 时已知。
\end{itemize}

\begin{proposition*}
$H^*_{N,\mathrm{cusp}}$ 的任何不可约分量都不是 $r$-可忽略的。所有分量都以正重数出现，且都是权 $0$ 纯的（$\Leftrightarrow$ Ramanujan--Petersson 猜想）。
\end{proposition*}

\par\medskip
\noindent\textbf{证明.}\

\begin{itemize}
\item 在引理的陈述中，可以把 $H^*_c(\Cht^{r,\ol{p}\le p}_N/a^{\ZZ})$ 换成 $IH^*_c(\ol{\Cht}^{r,\ol{p}\le p}_N/a^{\ZZ})$（这使纯性论证成为可能）。
\origpage{15}
\item 设 $H = H^*_{N,\mathrm{cusp}}(r-1)$。$H$ 的分量的权落在 $]-2, 2[$ 中（Jacquet 与 Shalika 的估计），并且是对称的。

设 $q'^*\sigma' \otimes q''^*\sigma''$ 是 $W_{(X-N)\times(X-N)}$ 的一个 $r$-可忽略的、权 $0$ 纯的表示。则 $\sigma' = \sigma_{\pi'}$、$\sigma'' = \sigma_{\pi''}$，其中 $\pi' \in \{\pi\}_{r'}$、$\pi'' \in \{\pi\}_{r''}$，$r', r'' < r$。
\item 对 $\pi \in \{\pi\}^N_r$，函数 $L(\pi \times \wt{\pi}', Z)$ 在区域 $|Z| < q^{-1+\varepsilon}$（$\varepsilon > 0$ 足够小）内没有极点，且级数
\[
\begin{aligned}
\frac{L'_{X-N}(\pi \times \wt{\pi}', Z)}{L_{X-N}(\pi \times \wt{\pi}', Z)}
&= \sum_{\infty \in (X-N)} \deg(\infty) \sum_{k \ge 1} \big(z_1(\pi_\infty)^{-k} + \cdots + z_r(\pi_\infty)^{-k}\big)\\
&\qquad \big(z_1(\pi'_\infty)^k + \cdots + z_{r'}(\pi'_\infty)^k\big)\, Z^{k\deg(\infty)-1}\\
&= \sum_{s \ge 1} Z^{s-1} \sum_{\substack{\infty \in |X-N|\\ \frac{s}{\deg(\infty)} = s' \in \mathbb{N}}} \deg(\infty)\\
&\qquad \big(z_1(\pi_\infty)^{-s'} + \cdots + z_r(\pi_\infty)^{-s'}\big)\big(z_1(\pi'_\infty)^{s'} + \cdots + z_{r'}(\pi'_\infty)^{s'}\big)
\end{aligned}
\]
在 $|Z| < q^{-1+\varepsilon}$ 内绝对收敛。
\item 类似地处理 $\dfrac{L'_{X-N}(\pi \times \wt{\pi}'', Z)}{L_{X-N}(\pi \times \wt{\pi}'', Z)}$。
\item 于是级数“乘积”（逐项相乘）
\[
\begin{aligned}
&\sum_{n \ge 1} Z^{n-1} \sum_{\substack{\infty, 0 \in |X-N|\\ \frac{s}{\deg(\infty)} = s' \in \mathbb{N},\ \frac{s}{\deg(0)} = u' \in \mathbb{N}}} \deg(\infty)\deg(0)\\
&\qquad \big(z_1(\pi_\infty)^{-s'} + \cdots + z_r(\pi_\infty)^{-s'}\big)\big(z_1(\pi_0)^{u'} + \cdots + z_r(\pi_0)^{u'}\big)\\
&\qquad \big(z_1(\pi'_\infty)^{s'} + \cdots + z_{r'}(\pi'_\infty)^{s'}\big)\big(z_1(\pi''_0)^{u'} + \cdots + z_{r''}(\pi''_0)^{u'}\big)
\end{aligned}
\]
对 $|Z| < q^{-2+2\varepsilon}$ 绝对收敛。
\item 由于 $\sigma' = \sigma_{\pi'}$、$\sigma'' = \sigma_{\pi''}$，级数
\[ \frac{L'_{(X-N)\times(X-N)}(H \otimes (\wt{\sigma}' \otimes \wt{\sigma}''), Z)}{L_{(X-N)\times(X-N)}(H \otimes (\wt{\sigma}' \otimes \wt{\sigma}''), Z)} \]
对 $|Z| < q^{-2+2\varepsilon}$ 绝对收敛。
\end{itemize}

$\Rightarrow$ 所以 $H$ 的权最大（$\geq 0$）的那一块不含任何 $r$-可忽略分量。

\origpage{16}

$\Rightarrow$ 若它有权 $> 0$，则权为 $1$，它会来自 $H^{2r-1}(\Cht^{r,\ol{p}\le p}/a^{\ZZ})$，并使系数出现 $-1$。

$\Rightarrow$ $H$ 是权 $0$ 纯的，且没有 $r$-可忽略分量。

\par\medskip
\noindent\textbf{以上讨论的结论：}\ 上同调的本质部分（$=$ 非 $r$-可忽略部分）集中在中间次数 $2r - 2$。它在

\[ H^{2r-2}_c(\Cht^{r,\ol{p}\le p}_N/a^{\ZZ}), \qquad IH^{2r-2}_c(\ol{\Cht}^{r,\ol{p}\le p}_N/a^{\ZZ}) \]

以及 $H^{2r-2}_c(\ol{\Cht}^{r,\ol{p}\le p'}_N/a^{\ZZ})$ 中都是一样的。

\par\medskip
\noindent\textbf{$\cH^r_N$（以及 $\Frob_\infty$、$\Frob_0$）的作用：}\ 设

\[ H^{2r-2}_c(\Cht^r_N/a^{\ZZ}) = \varinjlim_{p} H^{2r-2}_c(\Cht^{r,\ol{p}\le p}_N/a^{\ZZ}). \]

它带有以下作用：

\begin{itemize}
\item Hecke 代数 $\cH^r_N$；
\item $W_{(X-N)\times(X-N)}$；
\item 在 $\Frob_X \times \mathrm{Id}_X$ 与 $\mathrm{Id}_X \times \Frob_X$ 之上的两个“部分 Frobenius 自同态” $\Frob_\infty$、$\Frob_0$（且有 $\Frob_\infty \circ \Frob_0 = \Frob_0 \circ \Frob_\infty = \Frob$）。
\end{itemize}

\noindent$\bullet$ 这归结为

\[ \cH^r_N \times \big(W_{(X-N)\times(X-N)} \rtimes \ZZ\big) \]

的一个作用。

\begin{lemma*}
存在有限滤过
\[ 0 = H_0 \subsetneq H_1 \subsetneq \cdots \subsetneq H_m = H^{2r-2}_c(\Cht^r_N/a^{\ZZ}) \]
\emph{被 $\cH^r_N \times (W_{(X-N)\times(X-N)} \rtimes \ZZ)$ 保持，使得}
\begin{itemize}
\item \emph{$\forall i$，$H_{2i+1}/H_{2i}$ 只有 $r$-可忽略分量；}
\item \emph{$\forall i$，$H_{2i+2}/H_{2i+1}$ 只有本质分量。}
\end{itemize}
\end{lemma*}

\origpage{17}

\par\medskip
\noindent\textbf{证明.}\ 记 $H = H^{2r-2}_c(\Cht^r_N/a^{\ZZ})$。

\begin{itemize}
\item 令 $H_0 = 0$。
\item 对已定义的 $H_{2i}$，令：\\
$H_{2i+1}/H_{2i}$ = $H/H_{2i}$ 中所有有限维、$r$-可忽略的子表示的直和（$\Leftrightarrow$ 其不可约商都是 $r$-可忽略的）；\\
$H_{2i+2}/H_{2i+1}$ = $H/H_{2i+1}$ 中所有有限维、本质的子表示的直和。
\item 于是：
\begin{itemize}
\item 各 $H_i$ 被所有作用稳定；
\item 若 $i \ge 1$ 且 $H/H_i \neq 0$，则 $H_i \subsetneq H_{i+1}$；
\item 由于 $\bigoplus_i H_{2i+2}/H_{2i+1}$ 的维数有界，故 $i \gg 0$ 时 $H_i = H$。
\end{itemize}
\end{itemize}

$\Rightarrow$ 由此得到 $\cH^r_N \times (W_{(X-N)\times(X-N)} \rtimes \ZZ)$ 的一个半单表示 $H_{N,\mathrm{cusp}}$，其作为 $W_{(X-N)\times(X-N)}$ 表示的半单化正是 $H^*_{N,\mathrm{cusp}}$。

\begin{theorem*}
设 $f \in \cH^r_N$。对每个闭点 $x \in |X \times X|$，在其上 $\infty \neq 0 \in |X - N|$，并对 $\deg(x)$ 的每个倍数 $s = \deg(\infty)s' = \deg(0)u'$，有
\[
\begin{aligned}
\Tr_{H_{N,\mathrm{cusp}}}\big(f \times \Frob_x^{-\frac{s}{\deg(x)}}\big)
&= q^{(r-1)s} \sum_{\pi \in \{\pi\}^N_r} \Tr_\pi(f)\\
&\quad \big(z_1(\pi_\infty)^{-s'} + \cdots + z_r(\pi_\infty)^{-s'}\big)\big(z_1(\pi_0)^{u'} + \cdots + z_r(\pi_0)^{u'}\big).
\end{aligned}
\]
\end{theorem*}

\par\medskip
\noindent\textbf{证明概要.}\

\begin{itemize}
\item $f \in \cH^r_N$ 在 $H_{N,\mathrm{cusp}}$ 上的作用，与它在 $H^{2r-2}_c(\cX)$ 的本质部分上的作用相同，其中
\[
\cX = \ol{\Cht}^{r,\ol{p}\le p}/a^{\ZZ} \quad \text{若 } N = \emptyset, \qquad
\cX = \ol{\Cht}^{r,\ol{p}\le p'}_N/a^{\ZZ} \quad \text{若 } N \neq \emptyset.
\]
\item 回忆 $\cX$ 的边界是 DCN，其分层 $\cX_{\ul{r}}$ 由划分 $\ul{r} = (r_1 + \cdots + r_k = r)$ 编号，$k = |\ul{r}| \ge 2$。有一般的不动点公式：
\end{itemize}

\origpage{18}

\begin{proposition*}
存在边界分层 $\cX_{\ul{r}}$ 中的上同调对应 $\mathrm{cl}(f)_{\ul{r}}$，使得对每个 $x$ 与足够大的 $s$，
\[
\begin{aligned}
&\operatorname{Lef}_x(f \times \Frob^s, \Cht^{r,\ol{p}\le p}_N/a^{\ZZ}) = \sum_\nu (-1)^\nu \Tr_{H^\nu_c(\cX)}\big(f \times \Frob_x^{-\frac{s}{\deg(x)}}\big)\\
&\qquad + \sum_{\ul{r}} (-1)^{|\ul{r}|-1} \sum_\nu (-1)^\nu \Tr_{H^\nu_c(\cX_{\ul{r}})}\big(\mathrm{cl}(f)_{\ul{r}} \times \Frob_x^{-\frac{s}{\deg(x)}}\big).
\end{aligned}
\]
\end{proposition*}

\par\medskip
\noindent\textbf{所需的材料：}\ 这受过 Pink 的论证启发。

\begin{itemize}
\item 伴随公式。
\item 在 $\cX = \ol{\Cht}^{r,\ol{p}\le p}/a^{\ZZ}$ 或 $\ol{\Cht}^{r,\ol{p}\le p'}_N/a^{\ZZ}$ 中，Hecke 对应 $f \in \cH^r_N$ 稳定 chtoucas 的开集
\[ \cX_\emptyset = \Cht^{r,\ol{p}\le p}_N/a^{\ZZ} \]
 “在其不动点附近”。
\item 在 $\cX$ 不紧合的情形（$N \neq \emptyset$），Fujiwara 关于 Deligne 猜想的定理（用以消去不动点公式中无穷远处的项）。
\end{itemize}

\par\medskip
\noindent\textbf{定理证明的完成.}\

\begin{itemize}
\item 把计数公式与命题中的不动点公式结合起来，得到
\[
\begin{aligned}
&\sum_\nu (-1)^\nu \Tr_{H_c(\cX)}\big(f \times \Frob_x^{-s/\deg(x)}\big)\\
&\qquad + \sum_{\ul{r}} (-1)^{|\ul{r}|-1}\sum_\nu (-1)^\nu \Tr_{H_c(\cX_{\ul{r}})}\big(\mathrm{cl}(f)_{\ul{r}} \times \Frob_x^{-s/\deg(x)}\big)\\[2pt]
&= q^{(r-1)s} \sum_{\pi \in \{\pi\}^N_r} \Tr_\pi(f)\big(z_1(\pi_\infty)^{-s'} + \cdots + z_r(\pi_\infty)^{-s'}\big)\big(z_1(\pi_0)^{u'} + \cdots + z_r(\pi_0)^{u'}\big)\\[2pt]
&\quad + \text{其余项，其中出现的是各 } \GL_{r_1} \times \cdots \times \GL_{r_k} \text{ 的尖点自同调表示的 Hecke 特征值},\\
&\qquad\qquad r_1 + \cdots + r_k = r \qquad (k \ge 2).
\end{aligned}
\]
\origpage{19}
\item 最后，把公式两边的“本质部分”等同起来即得结论。
\end{itemize}

\par\medskip
\noindent\textbf{所需的材料：}

\begin{itemize}
\item 各 $H^\nu_c(\cX_{\ul{r}})$ 是 $r$-可忽略的；
\item 对应 $\pi' \mapsto \sigma_{\pi'}$ 在秩 $< r'$ 时已知 \quad $\Big\}$ 归纳假设；
\item 配对 $L$-函数（自守侧与 Galois 侧）的论证。
\end{itemize}

\begin{corollary*}
对每个 $\pi \in \{\pi\}^N_r$，存在 $W_{(X-N)\times(X-N)} \rtimes \ZZ$ 的 $\ell$-进表示 $H_\pi$，纯权为 $0$，使得对每个 $x \mapsto (\infty, 0)$ 与 $\deg(x)$ 的每个倍数 $s = \deg(\infty)s' = \deg(0)u'$，
\[ \Tr_{H_\pi}\big(\Frob_x^{-s/\deg(x)}\big) = \big(z_1(\pi_\infty)^{-s'} + \cdots + z_r(\pi_\infty)^{-s'}\big)\big(z_1(\pi_0)^{u'} + \cdots + z_r(\pi_0)^{u'}\big). \]
\end{corollary*}

\begin{theorem*}
对 $\pi \in \{\pi\}^N_r$，存在 $W_{X-N}$ 的一个纯权 $0$、$r$ 维不可约表示 $\sigma$，使得
\[ H_\pi \cong q'^*\sigma \otimes q''^*\check{\sigma} \qquad (q', q'' : X \times X \rightrightarrows X) \]
\emph{作为 $W_{(X-N)\times(X-N)}$ 的表示成立，并且 $\sigma = \sigma_\pi$。}
\end{theorem*}
